\chapter{Esercizi SQL}
\section{Totale Giugno 2026}
Si consideri il seguente schema relazionale:


\texttt{DIPENDENTE}(\underline{ID\_impiegato}, CF, cognome, nome, data\_assunzione)


\texttt{MANSIONE}(\underline{codice}, titolo, stipendio\_base)


\texttt{PREMIO\_PRODUZIONE}(\underline{dipendente}, mansione, data\_assegnazione, importo\_premio)

\subsection{Query 1}
Il titolo e lo stipendio base delle mansioni ordinate prima in ordine decrescente per stipendio base e poi per ordine alfabetico.

\begin{lstlisting}
CREATE VIEW AS info
SELECT titolo, stipendio_base
FROM MANSIONE;

(
    SELECT * FROM info
    ORDER BY stipendi_base DESC
)
UNION ALL
(
    SELECT * FROM info
    ORDER BY titolo ASC
)
\end{lstlisting}

\subsection{Query 2}
Il codice e il titolo delle mansioni che valgono il maggior stipendio base inserito a sistema.

\begin{lstlisting}
SELECT M.codice, M.titolo
FROM MANSIONE AS M
WHERE M.stipendio_base = (
    SELECT MAX(stipendio_base)
    FROM MANSIONE
);
\end{lstlisting}

\subsection{Query 3}
Il codice, il titolo della mansione e la media dell'importo del premio per ciascuna mansione presente nella tabella dei premi

\begin{lstlisting}
SELECT M.codice, M.titolo, AVG(P.importo_premio) AS media_premio
FROM MANSIONE AS M 
JOIN PREMIO_PRODUZIONE AS P ON M.codice = P.mansione
GROUP BY M.codice, M.titolo;
\end{lstlisting}

\subsection{Query 4}
L'ID impiegato, il nome e il cognome del dipendente che ha ricevuto il premio di produzione più basso in assoluto all'interno della mansione avente codice 'MAN01'.

\begin{lstlisting}
SELECT D.ID_impiegato, D.nome, D.cognome
FROM DIPENDENTE AS D JOIN PREMIO_PRODUZIONE AS P ON D.Id_impiegato = P.dipendente
WHERE P.importo_premio = (
    SELECT MIN(importo_premio) 
    FROM PREMIO_PRODUZIONE 
    WHERE mansione = 'MAN01'
)
\end{lstlisting}

\newpage
\section{Totale Giugno 2015}
Sia dato il seguente schema che descrive il database di un negozio di elettronica di comsumo

\begin{itemize}
    \item prodotto(\underline{codice}, nome, marca, prezzo)
    \item scontrino (\underline{progressivo}, data, ora, articolo, tesseraCliente)
    \item cliente (\underline{Tessera}, nome, cognome, cittaResidenza, puntiFedelta)
\end{itemize}
\subsection{Query 1}
Trovare il nome del prodotto più venduto tra le 11 e le 12 di qualunque giorno
\begin{lstlisting}
SELECT P.nome 
FROM Prodotto AS P
JOIN SCONTRINO AS S ON S.articolo = P.codice
GROUP BY P.codice, P.nome
HAVING COUNT(*) >= ALL (
    SELECT COUNT
    FROM SCONTRINO AS S2
    JOIN SCONTRINO AS S2 ON S.articolo = P.codice
    WHERE S2.ora BETWEEN '11:00' AND '12:00'
    GROUP BY S2.articolo
);
\end{lstlisting}


\subsection{Query 2}
Trovare il nome e il cognome dei clienti che pur avendo comprato un prodotto della "Apple" non hanno mai comprato un prodotto della "Samsung"

\begin{lstlisting}
SELECT C.nome, C.cognome
FROM CLIENTE AS C
WHERE EXISTS (
    SELECT P.marca
    FROM PRODOTTO AS P JOIN SCONTRINO AS S ON P.codice = S.articolo
    WHERE P.marca = "Apple" AND S.tesseraCliente = Cliente.tesseraCliente
)
AND
NOT EXISTS (
    SELECT P.marca
    FROM PRODOTTO AS P JOIN SCONTRINO AS S ON P.codice = S.articolo
    WHERE P.marca = "Samsung" AND S.tesseraCliente = Cliente.tesseraCliente
)
\end{lstlisting}

\newpage
\subsection{Query 3}
Trovare il nome e cognome dei clienti di Milano che hanno comprato contemporaneamente un prodotto "Philips" e un prodotto "Dell"

\begin{lstlisting}
SELECT C.nome, C.cognome
FROM CLIENTI AS C JOIN SCONTRINO AS S ON C.Tessera = S.tesseraCliente
WHERE C.cittaResidenza = 'Milano' AND EXISTS (
    SELECT P.marca
    FROM PRODOTTO AS P JOIN SCONTRINO AS S1 ON P.codice = S1.articolo
    WHERE P.marca = 'Philips' AND S1.tesseraCliente = C.Tessera
)
AND EXISTS (
    SELECT P.marca
    FROM PRODOTTO AS P JOIN SCONTRINO AS S1 ON P.codice = S1.articolo
    WHERE P.marca = 'Dell' AND S1.tesseraCliente = C.Tessera
)
\end{lstlisting}

\newpage
\section{2° Parziale 2016}
\begin{itemize}
    \item STUDENTE (\underline{matr}, nome, cognome, LuogoNascita) 
    \item ESAME (\underline{matr}, \underline{codCorso}, voto, data) 
    \item CORSO(\underline{codCorso}, nomeCorso, docente)
\end{itemize}

\subsection{Query 1}
Trovare la città in cui sono nati gli studenti che complessivamente hanno superato più esami rispetto agli studenti nati in altre città

\begin{lstlisting}
SELECT S.LuogoNascita
FROM STUDENTE AS S JOIN ESAME AS E ON S.matr = E.matr
GROUP BY S.LuogoNascita HAVING COUNT(*) >= ALL (
    SELECT COUNT(*)
    FROM STUDENTE AS S1 JOIN ESAME AS E1 ON S1.matr = E1.matr
    WHERE E1.voto >= 18
    GROUP BY S1.LuogoNascita
)
\end{lstlisting}

\subsection{Query 2}
Trovare nome cognome e matricola, degli studenti che hanno superato l'esame di analisi con voto superiore al voto ottenuto per l'esame di fisica

\begin{lstlisting}
SELECT S.nome, S.cognome, S.matr
FROM STUDENTE AS S JOIN ESAME AS E1 ON S.matr = E1.matr
JOIN CORSO AS C1 ON E1.codCorso = C1.codCorso
JOIN ESAME AS E2 ON S.matri = E2.matr
JOIN CORSO AS C2 ON E2.codCorso = C2.codCorso
WHERE C1.nomeCorso = 'Analisi' AND C2.nomeCorso 0 'Fisica'
AND E1.voto > E2.voto
\end{lstlisting}

\subsection{Query 3}
Trovare il nome dei corsi di cui studenti omonimi hanno superato l'esame
\begin{lstlisting}
SELECT C.nomeCorso 
FROM CORSO AS C 
JOIN ESAME AS E1 ON C.codCorso = E1.codCorso
JOIN ESAME AS E2 ON C.codCorso = E2.codCorso
JOIN STUDENTE AS S1 ON E,matr = S1.matr
JOIN STUDENTE AS S2 ON E.matr = S2.matr
WHERE S1.matr <> S2.matr 
AND  S1.nome = S2.nome AND S1.cognome = S2.cognome 
AND E1.codCorso = E2.codCorso
AND E1.voto >= 18 AND E2.voto >= 18
\end{lstlisting}

\newpage
\section{Totale Gennaio 2017}
La seguente base di dati descrive l'attività di una società di noleggio di autovetture che ha diverse sedi in alcune città italiane.

\begin{itemize}
    \item \textbf{SEDE}(\underline{NomeSede}, indirizzo, città, telefono, direttore)
    \item \textbf{AUTOVETTURA}(\underline{targa}, modello, AnnoImmatricolazione, numeroposti, colore)
    \item \textbf{CLIENTE}(\underline{CF}, nome, cognome, telefono, città)
    \item \textbf{CONTRATTO}(\underline{auto}, \underline{cliente}, \underline{DataNoleggio}, SedeNoleggio, DataConsegna, SedeConsegna)
\end{itemize}

\subsection{Query 1}
Trovare il numero di autovetture noleggiate in ogni città, da clienti con residenza in una città diversa dalla città in cui avviene il noleggio.
\begin{lstlisting}
SELECT S.citta ,COUNT(*) AS numero_noleggi
FROM CONTRATTO AS C
JOIN CLIENTE AS CE ON C.cliente = C.CF
JOIN SEDE AS S ON S.NomeSede = C.SedeNoleggio
WHERE S.citta <> CE.citta
GROUP BY S.citta
\end{lstlisting}

\subsection{Query 2}
Trovare i clienti (indicando CF, nome e cognome) che hanno noleggiato almeno due macchine diverse nella stessa sede di noleggio.

\begin{lstlisting}
SELECT C.CF, C.nome, C.cognome
FROM CLIENTE AS C
JOIN CONTRATTO AS CO1 ON C.CF = C01.cliente
JOIN CONTRATTO AS C02 ON C.CF = C02.cliente
WHERE C01.auto <> C02.auto
AND C01.SedeNoleggio = C02.SedeNoleggio
\end{lstlisting}

\subsection{Query 3}
Trovare la sede di noleggio, indicandone anche la città, che ha realizzato il maggior numero di contratti della durata di un giorno (cioè in cui la macchina sia stata riconsegnata lo stesso giorno in cui è stata noleggiata).

\begin{lstlisting}
SELECT S.NomeSede, S.citta
FROM SEDE AS S
JOIN CONTRATTO AS C ON S.NomeSede = C.SedeNoleggio
GROUP BY S.Nome, S.citta
HAVING COUNT(*) >= ALL (
    SELECT COUNT(*)p
    FROM CONTRATTO
    WHERE DataConsegna = SedeConsegna
    GROUP BY SedeNoleggio
)
\end{lstlisting}

\newpage
\section{Totale Luglio 2017}
La seguente base di dati descrive l'archivio di una società di charter che organizza e vende vacanze in barca. La società dispone di una flotta di barche ed organizza viaggi con diverse destinazioni.

\begin{itemize}
    \item \textbf{VIAGGIO} (\underline{idCliente}, \underline{dataPartenza}, \underline{nomeBarca}, Skipper, portoPartenza, portoArrivo, costo)
    \item \textbf{CLIENTE} (\underline{idCliente}, nome, cognome, nazione)
    \item \textbf{FLOTTA} (\underline{NomeBarca}, modello, lunghezza, posti, annoImmatricolazione)
    \item \textbf{PORTO} (\underline{Nome}, nazione, posti, telefono, canaleRadio)
\end{itemize}

\subsection{Query 1}
Trovare la nazione che è stata coinvolta nella maggior parte dei viaggi (come porto di arrivo o come porto di partenza).

\begin{lstlisting}
SELECT P.nazione
FROM PORTO AS P
JOIN VIAGGIO AS V ON V.portoPartenza = P.Nome OR V.portoArrivo = P.Nome
GROUP BY P.nazione
HAVING COUNT(*) >= ALL (
    SELECT COUNT(*)
    FROM PORTO AS P2
    JOIN VIAGGIO AS V2 ON V2.portoPartenza = p2.Nome OR V2.portoArrivo = p2.Nome
    GROUP BY P2.nazione
)
\end{lstlisting}

\subsection{Query 2}
Trovare il nome delle barche che non sono mai state affittate.
\begin{lstlisting}
SELECT F.NomeBarca
FROM FLOTTA AS F
WHERE F.NomeBarca NOT IN (
    SELECT V.nomeBarca
    FROM VIAGGIO
)
\end{lstlisting}
\subsection{Query 3}
Trovare per ogni porto il costo medio dei viaggi da lì partiti che non abbiano imbarcato clienti della nazionalità del porto stesso.

\begin{lstlisting}
SELECT P.nome, AVG(V.costo)
FROM PORTO AS P, VIAGGIO AS V
JOIN CLIENTE AS C ON V.idCliente = C.idCliente
WHERE C.nazione <> P.nazione
GRUOUP BY P.Nome
\end{lstlisting}

\section{Esercizio 1}
Data la seguente base di dati:
\begin{itemize}
    \item \textbf{STUDENTE}(\underline{matricola}, CF, cognome, nome, data\_nascita)
    \item \textbf{INSEGNAMENTO}(\underline{codice}, nome, crediti)
    \item \textbf{ESAME}(\underline{studente}, \underline{insegnamento}, data, voto, lode)
\end{itemize}

\subsection{Query 1}
Il codice e il nome degli insegnamenti che valgono il maggior numero di crediti
\begin{lstlisting}
SELECT I.codice, I.nome
FROM INSEGNAMENTO AS I
WHERE I.crediti = ALL (
    SELECT I2.crediti
    FROM INSEGNAMENTO AS I2
)
\end{lstlisting}

\subsection{Query 2}
Matricola e media voto degli studenti nati nel 2002
\begin{lstlisting}
SELECT S.matricola, AVG(E.voto)
FROM STUDENTE AS S, ESAME AS E
JOIN S ON E.studente = S.matricola
WHERE S.data_nascita = '2002'
GROUP BY S.matricola
\end{lstlisting}

\subsection{Query 3}
La matricola, il nome, il cognome dello studente che ha preso il massimo voto nell'esame dell'insegnamento di codice 'BDD01'
\begin{lstlisting}
SELECT S.matricola, S.nome, S.cognome
FROM STUDENTE AS S
JOIN ESAME AS E ON S.matricola = E.studente
JOIN INSEGNAMENTO AS I ON I.codice = E.insegnamento
WHERE E.insegnamento = 'BDD01' AND E.voto = (
    SELECT MAX(VOTO)
    FROM ESAME 
    WHERE E.insegnamento = 'BDD01'
)
\end{lstlisting}

\newpage
\section{Esercizio 2}
Data la seguente base di dati

\begin{itemize}
    \item \textbf{Aeroporto}(\underline{Città}, Nazione, NumPiste)
    \item \textbf{Volo}(\underline{IdVolo}, GiornoSett, CittàPart, OraPart, CittàArr, OraArr, Aereo) 
    \item \textbf{Aereo}(\underline{TipoAereo}, NumPasseggeri, QtaMerci)
\end{itemize}

\subsection{Query 1}
Selezionare tutti gli ID dei voli che utilizzano l'AIRBUS 400

\begin{lstlisting}
SELECT V.IdVolo
FROM VOLO AS V
WHERE V.Aereo = 'AIRBUS 400'
\end{lstlisting}

\subsection{Query 2}
Selezionare gli aereoporti che hanno il magior numero di piste
\begin{lstlisting}
SELECT Citta
FROM Aereoporto AS A
WHERE A.NumPiste = (
    SELECT MAX(A1.NumPiste)
    FROM Aeroporto AS A1
)
\end{lstlisting}

\subsection{Query 3}
Selezionare le tipologie di aerei che hanno viaggiato in italia

\begin{lstlisting}
SELECT DISTINCT V.Aereo
FROM Volo AS V
JOIN Aereoporto AS AR1 ON V.CittaPart = AR1.Citta 
JOIN Aereoporto AS AR2 ON V.CittaArr = AR2.Citta
Where AR1.Nazione = 'Italia' OR AR2.Nazione = 'Italia'
\end{lstlisting}

\subsection{Query 4}
Selezionare le tipologie di aerei che hanno viaggiato maggiormente
\begin{lstlisting}
    
\end{lstlisting}

\subsection{Query 5}
Selezionare le tipologie di aerei che hanno effettuato più di 100 voli

\begin{lstlisting}
SELECT V.Aereo
FROM VOLO AS V
GROUO BY V.Aereo
HAVING COUNT(*) > 100
\end{lstlisting}

\newpage
\section{Esercizio 3}
Data la seguente base di dati che descrive un campionato invernale di sci maschile

\begin{itemize}
    \item \textbf{LUOGOGARA}(nome, altitudine, nazione) 
    \item \textbf{GARA} (specialità, luogo, dataGara, lunghezza, dislivello, premiovincitore) 
    \item \textbf{POSIZIONEGARA} (idatleta, dataGara, numeroPettorina, specialità, luogo, posizionearrivo) 
    \item \textbf{ATLETA} (idatleta, nome, cognome, nazione, datadinascita)
\end{itemize}

\subsection{Query 1}
Trovare per ogni nazione l'atleta che ha vinto il maggior numero di gare

\begin{lstlisting}
SELECT A.idatleta
FROM ATLETA AS A
JOIN POSIZIONEGARA AS P ON P.idatleta = A.idatleta
WHERE P.posizionearrivo = 1
GROUP BY A.nazione, A.idatleta
HAVING COUNT(*) >= ALL (
    SELECT COUNT(*)
    FROM ATLETA AS A2 JOIN POSIZIONEGARA AS P2 ON A2.idatleta = P2.idatleta
    WHERE posizionearrivo = 1 AND A.nazione = A2.nazione
    GROUP BY A2.idatleta
)
\end{lstlisting}

\subsection{Query 2}
Trovare le nazioni che hanno almeno un'atleta che ha vinto almeno una gara in patria

\begin{lstlisting}
SELECT DISTINCT A.nazione
FROM ATLETA AS A
JOIN POSIZIONEGARA AS P ON A.idatleta = P.idatleta
JOIN LUOGOGARA AS L ON L.nome = P.luogo
WHERE P.posizionearrivo = 1 AND L.nazione = A.nazione
\end{lstlisting}

\newpage
\section{Esercizio 4}
\begin{itemize}
    \item \textbf{CLIENTE}(\underline{CodiceFiscale}, Nome, Cognome, Nazionalità, Telefono, e-mail) 
    \item \textbf{STANZA}(\underline{CodiceStanza}, \underline{Albergo}, Tipo, MetriQuadri, Locali, Piano) 
    \item \textbf{SOGGIORNO}(\underline{CodiceStanza}, \underline{Albergo}, \underline{DataInizio}, DataFine, Agente, CFCliente, CostoGiornaliero) 
    \item \textbf{ALBERGO}(\underline{CodiceAlbergo}, NumeroDiStelle, Indirizzo, Città, Stato) 
    \item \textbf{AGENTE}(\underline{CodiceAgente}, Nome, Cognome, DataDiNascita)
\end{itemize}

\subsection{Query 1}
Trovare, specificando nome, cognome e codice fiscale, i clienti omonimi che abbiano soggiornato in una qualunque stanza di albergo nella medesima città

\begin{lstlisting}
SELECT C1.Nome, C1.Cognome, C1.CodiceFiscale
FROM CLIENTE AS C1
JOIN SOGGIORNO AS S1 ON C1.CodiceFiscale = S1.CFCLIENTE
JOIN ALBERGO AS AL1 ON S1.Albergo = AL1.CodiceAlbero
WHERE EXISTS (
    SELECT * 
    FROM CLIENTE AS C2
    JOIN SOGGIORNO AS S2 ON C2.CodiceFiscale = S2.CFCliente
    JOIN ALBERGO A2 ON S2.Albergo = A2.CodiceAlbergo
    WHERE C1.Nome = C2.nome AND C2.Cognome = C1.Cognome AND C1.CodiceFiscale <> C2.CodiceFiscale AND A1.Citta = A2.Citta
)
\end{lstlisting}

\subsection{Query 2}
Trovare l'agente (codice, nome, cognome) che ha stipulato più soggiorni per l'anno 2009 in alberghi con più di 3 stelle

\begin{lstlisting}
SELECT AG.CodiceAgente, AG.Nome, AG.Cognome
FROM AGENTE AS AG
JOIN SOGGIORNO AS S ON S.Agente = AG.CodiceAgente
JOIN ALBERO AS AL ON AG.Albergo = AL.CodiceALbergo
WHERE S.DataInizio >= 1-1-2009 AND S.DataInizio < 1-1-2010 AND AL.NumeroDiStelle > 3
GROUP BY AG.CodiceAgente, AG.Nome, AG.Cognome
HAVING COUNT(*) >= ALL
\end{lstlisting}

\newpage
\section{Esercizio 5}
\begin{itemize}
    \item \textsc{Palazzetto}(\underline{Nome}, Città, Capienza)
    \item \textsc{Incontro}(\underline{NomePalazzetto, Data, Ora}, Squadra1, Squadra2)
    \item \textsc{Nazionale}(\underline{Nazione}, Continente, Livello).
\end{itemize}

\subsection{Query 1}
Determinare i nomi dei palazzetti in cui non gioca nessuna nazionale asiatica
\begin{lstlisting}
SELECT P.Nome
FROM PALAZZETTO AS P
WHERE P.Nome NOT IN (
    SELECT I.NomePalazzetto
    FROM INCONTRO AS I
    JOIN NAZIONALE AS N1 ON I.Squadra1 = N1.Nazione
    JOIN NAZIONALE AS N2 ON I.Squadra2 = N2.Nazione
    WHERE N1.Continente = 'Asia' OR N2.Continente = 'Asia'
)
\end{lstlisting}


\subsection{Query 2}
Determinare la capienza complessiva dei palazzetti in cui si giocano partite di nazionali africane (ai fini della valutazione della capienza complessiva, si sommino le capienza associate a ciascuna gara, anche se più gare si svolgono nello stesso palazzetto)

\begin{lstlisting}
-- Subquery = Selezionare i palazzetti dove giocano le nazionali afrincane
-- OuterQuery = Sommora le capienze
SELECT SUM(P.Capienza)
FROM PALAZZETTO AS P
JOIN INCONTRO AS I ON P.Nome = I.NomePalazzetto
JOIN NAZIONALE AS N1 ON I.Squadra1 = N1.Nazione
JOIN NAZIONALE AS N2 ON I.Squadra2 = N2.Nazione
WHERE N1.Continente = 'Africa' OR N2.Continente = 'Africa'
\end{lstlisting}

\subsection{Query 3}
Individuare la città (le città se più di una soddisfa le condizioni dell'interrogazione) in cui si trova il palazzetto in cui la squadra olandese gioca il maggior numero di partite
\begin{lstlisting}
--SubQuery trova il palazzetto dove l'olanda gioca piu partite
--OuterQuery Trova il nome della citta di quel palazzetto
SELECT P.citta
FROM PALAZZETTO AS P
WHERE COUNT(P.Nome) >= ALL (
    SELECT P1.Nome
    FROM PALAZZETTO AS P1
    JOIN INCONTRO AS I ON I.NomePalazzetto = P1.Nome
    JOIN NAZIONE AS N1 ON I.Squadra1 = N1.Nazione
    JOIN NAZIONE AS N2 ON I.Squadra2 = N2.Nazione
    WHERE N1.Nazione = 'Olanda' OR N2.Nazione = 'Olanda'
)
\end{lstlisting}


\subsection{Query 4}
Determinare le squadre che incontrano solo squadre dello stesso livello
\begin{lstlisting}
SELECT N.Nazione
FROM NAZIONALE AS N
WHERE NOT EXISTS (
    SELECT *
    FROM INCONTRO AS I, NAZIONALE AS N2
    JOIN I ON I.Squadra1 = N.Nazione
    JOIN I ON I.Squadra = N2.Nazione
    WHERE N.Livello <> N2.Livello
)
AND
NOT EXISTS (
    SELECT *
    FROM INCONTRO AS I, NAZIONALE AS N2
    JOIN I ON I.Squadra1 = N2.Nazione
    JOIN I ON I.Squadra = N.Nazione
    WHERE N2.Livello <> N.Livello
)
\end{lstlisting}

\newpage
\section{Esercizio 6}
Si consideri la seguente Base di Dati:

\begin{itemize}
    \item DISCO(\underline{Nroserie}, TitoloAlbum, Anno, Prezzo)
    \item CONTIENE(\underline{NroSerieDisco, CodiceReg}, NroProgr)
    \item ESECUZIONE(\underline{CodiceReg}, TitoloCanz, Anno)
    \item AUTORE(\underline{Nome, TitoloCanzone})
    \item CANTANTE(\underline{NomeCantante, CodiceReg})
\end{itemize}

\subsection{Query 1}
La canzone che è stata eseguita almeno due volte
\begin{lstlisting}
SELECT E.TitoloCanz, COUNT(*) AS NumeroEsecuzioni
FROM ESECUZIONE AS E
GROUP BY E.TitoloCanz
HAVING COUNT (*) >= 2
\end{lstlisting}

\subsection{Query 2}
Gli autori di canzoni eseguite tra il 2000 e il 2010
\begin{lstlisting}
SELECT A.Nome
FROM AUTORE AS A
JOIN ESECUZIONE AS E ON A.TitoloCanzone = E.TitoloCanz
WHERE E.anno >= 2000 AND E.anno <= 2010
\end{lstlisting}

\subsection{Query 3}
Gli autori delle canzoni più eseguite
\begin{lstlisting}
SELECT A.Nome ,A.TitoloCanzone, COUNT(*)
FROM AUTORE AS A 
JOIN ESECUZIONE AS E ON E.TitoloCanz = A.TitoloCanzone
GROUP BY A.Nome, A.TitoloCanzone
HAVING COUNT(*) >= ALL (
    SELECT COUNT(*)
    FROM ESECUZIONE
    GROUP BY TitoloCanz
)
\end{lstlisting}

\newpage
\section{Esercizio 7}
\noindent
Data la seguente base di dati

\begin{itemize}
    \item Concorso(\underline{Codice}, NumeroProve, Bando)
    \item Bando(\underline{Codice}, DataPubblicazione, DataScadenza)
    \item Candidato(\underline{CF}, Cognome, Nome, Indirizzo, Telefono)
    \item Commissario(\underline{CF}, Cognome, Nome, NumeroTelefono)
    \item Prova(\underline{Concorso, Data}, PunteggioMassimo)
    \item Svolgimento(\underline{DataProva, ConcorsoProva, Candidato}, Punteggio)
    \item Partecipazione(\underline{Concorso, Candidato}, PunteggioTotale, PosizioneInGraduatoria)
    \item Risposta(\underline{Bando, Candidato})
    \item Commissione(\underline{Concorso, Commissario}, Ruolo)
\end{itemize}

\subsection{Query 1}
Selezionare tutti i concorsi del 2018
\begin{lstlisting}
SELECT C.*
FROM Concorso AS C
JOIN Bando AS B ON C.Bando = B.Codice
WHERE (B.DataPubblicazione = '2018-12-1') OR (B.DataScadenza = '2018-12-31')
\end{lstlisting}

\subsection{Query 2}
Selezionare tutti i candidati che hanno fatto concorsi con bando scaduto nel 2018
\begin{lstlisting}
SELECT C.*
FROM CANDIDATO AS C
JOIN Partecipazione AS P ON C.CF = P.Candidato
JOIN CONCORSO AS CON ON P.Concorso = CON.Codice
JOIN BANDO AS B ON CON.Bando = B.Codice
WHERE YEAR(B.DataScadenza) = 2018
\end{lstlisting}

\subsection{Query 3}
Trovare i commissari dei concorsi con numero di prove pari a 2
\begin{lstlisting}
SELECT CO.*
FROM COMMISSARIO AS CO
JOIN Commissione AS COM ON CO.CF = COM.Commissario
JOIN CONCORSO AS CON ON COM.Concorso = CON.Codice
WHERE CON.NumeroProve = 2
\end{lstlisting}

\newpage
\section{Esercizio 8}
\noindent
Data la seguente base di dati

\begin{itemize}
    \item Lavoratore(\underline{CF}, Nome, Cognome, Età)
    \item Città(\underline{Nome}, NumeroAbitanti)
    \item HaAbitato(\underline{CFPersona, Città}, DaAnno, AAnno)
    \item HaLavorato(\underline{CFPersona, PIVADitta, Città}, DaAnno, AAnno)
    \item Ditta(\underline{PIVA}, Nome, NumeroImpiegati, CapitaleSociale)
    \item HaAvutoSedeIn(\underline{PIVADitta, NomeCittà}, DaAnno, AAnno)
\end{itemize}

\subsection{Query 1}
Trovare i lavoratori delle aziende milanesi
\begin{lstlisting}
SELECT L.*
FROM LAVORATORE AS L
JOIN HaLavorato AS HL ON HL.CFPersona = L.CF
WHERE HL.Citta = 'Milano'
\end{lstlisting}

\subsection{Query 2}
Trovare le partite IVA delle ditte che hanno (o hanno avuto) sedi in tutte le città;
\begin{lstlisting}
-- Inserisci qui la query 2
\end{lstlisting}

\subsection{Query 3}
Per ogni ditta trovare il numero degli impiegati che vi hanno lavorato, risiedendo nella stessa città dove aveva sede la ditta
\begin{lstlisting}
\end{lstlisting}

\subsection{Query 4}
Trovare le città che hanno o hanno avuto il maggior numero di ditte
\begin{lstlisting}
-- Inserisci qui la query 4
\end{lstlisting}

\subsection{Query 5}
Trovare le città che hanno avuto il maggior numero di ditte con sede nel 2019
\begin{lstlisting}
-- Inserisci qui la query 5
\end{lstlisting}

\newpage
\section{Esercizio 9}
\begin{itemize}
    \item CINEMA(\underline{codice}, nome, indirizzo, codice\_comune)
    \item COMUNE(\underline{codice}, nome, provincia, regione, numero\_abitanti)
    \item FILM(\underline{codice}, titolo, anno, nazionalita, tecnica\_animazione)
    \item PROIETTATO\_IN(\underline{codice\_film, codice\_sala, data\_ora}, lingua)
    \item SALA(\underline{codice}, nome, capienza, abilitata\_3d, codice\_cinema)
\end{itemize}

\subsection{Query 1}
Selezionare tutti i film di tecnica animazione «computer graphics» proiettati nel 2016
\begin{lstlisting}
SELECT DISTINCT F.*
FROM FILM AS F
JOIN PROIETTATO_IN AS P ON F.codice = P.codice_film
WHERE F.tecninca_animazione = 'computer graphics' AND YEAR(P.data_ora) = 2016
\end{lstlisting}

\subsection{Query 2}
Selezionare la capienza media delle sale dei cinema che hanno più di una sala
\begin{lstlisting}
SELECT S.codice_cinema ,AVG(S.capienza)
FROM SALA AS S 
GROUP BY S.codice_cinema
HAVING COUNT(*) >= 2
\end{lstlisting}

\subsection{Query 3}
Selezionare i cinema e nome della città che hanno il maggior numero di posti a vedere.
\begin{lstlisting}
SELECT C.codice, CO.nome, SUM(S.capienza)
FROM CINEMA AS C
JOIN C ON C.codice_comune = CO.codice
JOIN S ON S.codice_cinema = C.codice
GROUP BY C.codice, CO.nome
HAVING SUM(S.capienza) >= ALL (
    SELECT SUM(S2.capienza)
    FROM SALA AS S2
    GROUP BY S2.codice_cinema
)
\end{lstlisting}